\scanpage{26}
% L1–4
\uline{Exercice} Soit $B\xrightarrow{p}I$ une application \inserted{\unsure{surjective|surjection}} \deleted{telle que
$\forall i\in I$, $\operatorname{card}p^{-1}(i)\geq2$.} Soit $\mathcal P$ l'ens des sections
\unsure{partiellt|partiellement|partielles} définies de $B$ sur (des parties de) $I$,
\deleted{\unsure{muni de cette}} ordonné par la relation de prolongement.

% L5–8
\deleted{Montrer que les} \deleted{\unsure{et l'ens}} Ainsi l'ens des él. maximaux
de $\mathcal P$ s'identifie à $C=\prod_{i\in I}B_i$ $(B_i=p^{-1}(\unsure{\{i\}|i}))$
L'ens des éléments \deleted{\unsure{minimi|minim.}} \inserted{sous-minimaux} s'identifie à
$B$. Pour tt $b\in B$, soit \deleted{\unsure{F_b}} $C_b\subset C$ l'ens
% L9–13
des él. maximaux \deleted{tels que} qui majorent $b$,
i.e. l'ens des sections $\sigma$ de $B$ sur $I$ telles
que $\sigma(p(b))=b$. Soient $b,b'\in B$, \inserted{$b\ne b'$,} alors
on a
\[
 C_b\cap C_{b'}=\varnothing\quad\Longleftrightarrow\quad p(b)=p(b').
\]
Donc \inserted{si $b,b'\in B$, on a}
\[
 (p(b)=p(b'))\quad\Longleftrightarrow\quad(b=b'\ \text{ou}\ C_b\cap C_{b'}=\varnothing)
\]

% L14–17
Cela nous permet de reconstruire $B$ et
la relation d'équivalence \inserted{$R$} définie par $p$ \deleted{\unsure{en}}
\deleted{\unsure{la}} (donc $I=B/R$ et $p$), à isom près, en
termes de $\mathcal P$.

% L18–21
\uline{Cor.} La cat. des $B\xrightarrow{p}I$ ($p$ surjectif)
(pour les iso) équivalente à une sous-catégorie
pleine de celle des ens. ordonnés
en associant à tt $p$ \inserted{le} $\mathcal P(p)$ ci-dessus
% archive sheet number: 25
